\section{Данные}

\begin{frame}[c]{Данные: сколько и откуда}
  \small Чем сложнее задача, тем больше данных нужно. Они должны быть и многочисленными, и качественными.

  \smallskip
  \begin{columns}[T]
    \begin{column}{0.47\textwidth}
      \begin{block}{Открытые датасеты}
        Их много, но чаще для простых задач, например классификации картинок.
      \end{block}
      \begin{block}{Краудсорсинг}
        Люди размечают данные. Волонтёры без оплаты --- citizen science.
      \end{block}
    \end{column}
    \begin{column}{0.47\textwidth}
      \begin{block}{Покупка}
        Иногда датасет можно купить. Для некоторых задач данных нет нигде.
      \end{block}
      \begin{block}{Неразмеченные данные}
        Self-supervised learning и обучение представлений: используем данные без ответов.
      \end{block}
    \end{column}
  \end{columns}
\end{frame}

\begin{frame}[c]{Типы признаков}
  \small Признаки --- свойства объекта, по которым модель его описывает. Удобнее всего таблица: строки --- объекты, столбцы --- признаки.

  \smallskip
  \begin{columns}[T]
    \begin{column}{0.47\textwidth}
      \begin{block}{Численный}
        Рост, доход. Иногда отдельно выделяют вещественные и целочисленные.
      \end{block}
      \begin{block}{Бинарный}
        Два значения: да или нет, 0 или 1. Можно считать и числом, и категорией.
      \end{block}
    \end{column}
    \begin{column}{0.47\textwidth}
      \begin{block}{Категориальный}
        Значение из набора вариантов: профессия, день недели.
      \end{block}
      \begin{block}{Ординальный}
        Категории с порядком: класс опасности вещества от 1 до 4, год обучения.
      \end{block}
    \end{column}
  \end{columns}
\end{frame}

\begin{frame}[c]{Какого типа признак?}
  \lead{Вопрос на подумать.}
  \small Численный $\cdot$ Категориальный $\cdot$ Бинарный $\cdot$ Ординальный

  \smallskip
  \begin{columns}[T]
    \begin{column}{0.47\textwidth}
      \begin{numrows}
        1 & Рост человека\\
        2 & Профессия\\
        3 & Оплатил ли клиент заказ\\
      \end{numrows}\par
    \end{column}
    \begin{column}{0.47\textwidth}
      \begin{numrows}
        4 & Класс энергопотребления (от A до G)\\
        5 & День недели\\
        6 & Год обучения студента\\
      \end{numrows}\par
    \end{column}
  \end{columns}
\end{frame}

\begin{frame}[c]{Ответ: типы признаков}
  \small
  \zebra{2}
  \zebra{1}
  \begin{tabular}{>{\raggedright\arraybackslash}p{0.444\textwidth}l}
    \rowcolor{acDark} \textcolor{white}{\textbf{Признак}} & \textcolor{white}{\textbf{Тип}}\\
    Рост человека & \textcolor{acMid}{\textbf{Численный}}\\
    Профессия & \textcolor{acMid}{\textbf{Категориальный}}\\
    Оплатил ли заказ & \textcolor{acMid}{\textbf{Бинарный}}\\
    Класс энергопотребления, A--G & \textcolor{acMid}{\textbf{Ординальный}}\\
    День недели & \textcolor{acMid}{\textbf{Категориальный}}\\
    Год обучения & \textcolor{acMid}{\textbf{Ординальный}}\\
  \end{tabular}
\end{frame}

\begin{frame}[c]{Что не так с этими данными?}
  \lead{Найдите проблемы.}

  \medskip
  \large
  \zebra{2}\renewcommand{\arraystretch}{1.7}\setlength{\tabcolsep}{12pt}
  \begin{tabular}{llrl}
    \rowcolor{acDark} \textcolor{white}{\textbf{Клиент}} & \textcolor{white}{\textbf{Возраст}} & \textcolor{white}{\textbf{Доход, руб.}} & \textcolor{white}{\textbf{Вернул кредит}}\\
    Анна    & 34  & 90\,000    & да\\
    Борис   & N/A & 70\,000    & нет\\
    Виктор  & 140 & 85\,000    & да\\
    Галина  & 29  & 2\,400\,000 & да\\
    Дмитрий & 41  & 65\,000    & да\\
  \end{tabular}

  \bigskip
  \textbf{Часть проблем в таблице не видна вовсе. Какие это?}
\end{frame}

\begin{frame}[c]{Ответ: четыре типа проблем в данных}
  \begin{columns}[T]
    \begin{column}{0.47\textwidth}
      \begin{block}{Пропуски}
        Возраст Бориса неизвестен. Удалить, заполнить или оставить как признак: пропуск сам может нести информацию.
      \end{block}
      \begin{block}{Ошибки разметки}
        Разметчики ошибаются: бывают недобросовестные, а задача может быть сложной.
      \end{block}
    \end{column}
    \begin{column}{0.47\textwidth}
      \begin{block}{Выбросы}
        140 лет и доход 2,4 млн --- ошибка сбора или реальная аномалия. Важные выбросы обрабатывают отдельно.
      \end{block}
      \begin{block}{Data drift}
        Данные меняются со временем: формат или распределение. Нужен мониторинг.
      \end{block}
    \end{column}
  \end{columns}
\end{frame}
